% kernel_loewner08.tex -- round 496.
% PRIME.RDAGGER.KERNEL_LOEWNER.02.
% NO_GO(kernel-Loewner-compact-tail@L=0.8).
% Companion: verify_kernel_loewner08.py.
% Probe: experiments/tfpt-discovery/kernel_loewner08_probe.py.
% Method no-go only.  NO RH CLAIM.
\documentclass[11pt]{article}

\usepackage[a4paper,margin=2.45cm]{geometry}
\usepackage{amsmath,amssymb,amsthm}
\usepackage{booktabs,array,microtype,xcolor}
\usepackage[T1]{fontenc}
\usepackage[colorlinks=true,linkcolor=blue!50!black,urlcolor=blue!50!black]{hyperref}

\newtheorem{proposition}{Proposition}
\newtheorem{remark}{Remark}
\newcommand{\QW}{Q_{\mathrm W}}
\newcommand{\nogo}{\textcolor{red!65!black}{NO\_GO}}
\setlength{\emergencystretch}{2em}

\title{\bfseries The compact-tail kernel/Loewner method at $L=0.8$\\[2mm]
\large a sealed method no-go beyond the prime-free zone\\[2mm]
\large PRIME.RDAGGER.KERNEL\_LOEWNER.02 --- round 496}
\author{Stefan Hamann}
\date{August 31, 2026}

\begin{document}
\maketitle

\begin{abstract}\noindent
The round-492 three-gate contract is tested at $L=0.8$, where the prime
powers $2,3,4$ enter.  Gate G1 passes: the $x$-space identity, its
factor-two prime normalization, the zero-extension convention, and the
r471/r477 calibration stops agree.  Gates G2 and G3 give an honest
method no-go.  The best tested booked cutoff already has negative
$101$-section margin, while its full three-piece Hilbert--Schmidt
charge is $27.3672$.  More fundamentally, the prime translations are
bounded but noncompact; their measured Legendre off-block does not
vanish.  Therefore the compact-kernel completion used at $L=0.3$
cannot treat the prime section as the whole operator.  This excludes
one method only.  It proves neither negativity nor failure of Weil
positivity at $L=0.8$, and it is not a claim for or against RH.
\end{abstract}

\medskip\noindent
\textbf{Verdict.}\enspace
\texttt{NO\_GO(kernel-Loewner-compact-tail@L=0.8)}.

\smallskip\noindent
Probe SPEC\_SHA:
\nolinkurl{63f590d43c9ed63e8419c99ddd6b9110e58a7f419d6445de71932aa4575cba97}.

\begin{center}
\fbox{\parbox{0.95\linewidth}{%
\textbf{Claim boundary.}
This note certifies a failure of the prescribed compact-tail proof
budget, not a counterexample and not an anti-positivity theorem.
The positive r471 and r477 finite-subspace calibrations remain valid.
No Galerkin value is promoted to an operator value.  No RH claim.}}
\end{center}

\section{Exact prime normalization and the first balance}

For real $h$ supported in $[-L,L]$, extend $h$ by zero to
$L^2(\mathbb R)$ and put $g=h*\widetilde h$.  The r471 classical
normalization is
\[
 P(h)=2\sum_{n\in\{2,3,4\}}
       \frac{\Lambda(n)}{\sqrt n}\,g(\log n).
\]
Thus, with $a_n=\Lambda(n)/\sqrt n$,
\[
 -P(h)=-2\sum_n a_n g(0)
       +2\sum_n a_n\,
          \langle h,(I-\operatorname{Re}\tau_{\log n})h\rangle .
\]
The second sum is positive.  The first is the full prime $g(0)$
load.  The factor two is essential: omitting it misses both r471 and
r477.

At $L=0.8$ the interval computation gives
\begin{align*}
 c_L&=3.5598735700355446228812256099,\\
 (a_2,a_3,a_4)
  &=(0.4901290717342736,\ 0.6342841005975640,\
     0.3465735902799727),\\
 2\sum_n a_n&=2.9419735252236204555039093902.
\end{align*}
For the r494 cutoff $[0.01,0.03]$,
$\kappa_w=5.065476699202523$ and the scalar balance
$\kappa_w-c_L-2\sum a_n=-1.4363703961$ is already negative.
Only substantially harder cutoffs make that scalar line positive.

\section{Gate G1: identity passes}

Five compactly supported polynomial test functions were evaluated
independently in the $x$-space form and in the defining frequency
form at \texttt{mp.dps=50}.  The frequency multiplier included
\[
 -2\sum_{n\in\{2,3,4\}}a_n\cos(t\log n).
\]
The maximum normalized discrepancy was $5.499\cdot10^{-14}$.
After this gate, no frequency quantity was used again.

The support convention was tested by the r495-A2 rational staircase
battery: all six identities
\[
 2\{g(0)-g(a)\}=\|h-\tau_a h\|_{L^2(\mathbb R)}^2
\]
held exactly over $\mathbb Q$, including the boundary strips omitted
by an interval-only norm.

\begin{center}
\small
\begin{tabular}{@{}lll@{}}
\toprule
Calibration stop & Enclosure/value & Result\\
\midrule
r471, $L=0.8$, $n=24$
 & hint $2.64507571918113\cdot10^{-4}$,
   $\mu_{\rm lo}=1.32253673213734\cdot10^{-4}$ & pass\\
r477 constant mode
 & $r(1)\in[0.0802740548742,0.0802740551870]$ & pass\\
r477 $(1,\text{first Dirichlet})$
 & Schur $\in[0.0741382774823,0.0750735086403]$ & pass\\
\bottomrule
\end{tabular}
\end{center}

\section{Gate G2: the booked cutoff budget fails}

Let $w$ be zero on $[0,x_0]$, one on $[x_1,1.6]$, and the quintic
$C^2$ smoothstep between them.  Every row below includes
$c_L$, the full factor-two prime $g(0)$ mass, and the omitted
archimedean integral through the Loewner minorant.

\begin{center}
\small
\begin{tabular}{@{}ccrrr@{}}
\toprule
$x_0$ & $x_1$ & $\kappa_w$ & scalar balance
 & $\lambda_{\min}(P_{101}Q_wP_{101})$\\
\midrule
0.001&0.003&7.3770530633&$+8.752060\cdot10^{-1}$&$-1.6106733\cdot10^{-4}$\\
0.002&0.006&6.6829061429&$+1.810590\cdot10^{-1}$&$-8.1766226\cdot10^{-4}$\\
0.003&0.010&6.1992956339&$-3.025515\cdot10^{-1}$&$-2.2841017\cdot10^{-3}$\\
0.010&0.030&5.0654766992&$-1.436370\cdot10^{0}$&$-3.0322577\cdot10^{-2}$\\
\bottomrule
\end{tabular}
\end{center}

The best tested row already has negative finite-section margin.
By Cauchy interlacing, every containing Legendre section through
$n=800$ has minimum eigenvalue no larger than this value.  Hence
there is no positive finite margin to which a remainder could be
compared.

\section{Gate G3: compact-tail completion fails}

For
\[
 T_w(y,z)=\frac12w(|y-z|)
 \frac{e^{|y-z|/2}}{\sinh|y-z|},
\]
the exact Hilbert--Schmidt norm is
\[
 \|T_w\|_{\rm HS}^2
 =\frac12\int_0^{1.6}(1.6-x)
   \left(w(x)\frac{e^{x/2}}{\sinh x}\right)^2\,dx .
\]
For the best G2 row,
\[
 \|T_w\|_{\rm HS}=19.217696275,\qquad
 r_{101}=9.122410180,\qquad
 3r_{101}=27.367230540.
\]
The full $r+2r$ charge is used; neither $r/3$ nor the central block
alone is admitted.

There is a second, structural obstruction.  A compressed translation
has a delta-line kernel.  It is bounded, but it is not compact or
Hilbert--Schmidt.  Exact Gauss integration of the prime translation
matrix gives
\[
 \left\|P_{[101,301)}
   \left(2\sum_n a_n\operatorname{Re}R_{\log n}\right)P_{101}
 \right\|=0.939496613,
\]
with Frobenius norm $3.807582717$.  Consequently its finite Legendre
section cannot be substituted for the operator and cannot be absorbed
into the r494 Hilbert--Schmidt tail.  Doing so would be precisely the
forbidden Galerkin-as-operator pad.

\begin{proposition}[Round-496 method no-go]
Under the round-492/r494 three-gate contract, the compact-kernel
finite-section method does not prove
\[
 \lambda_*(0.8)\ge c
\]
for any $c>0$.  Gate G1 passes, but G2 has no positive certified
finite margin and G3 cannot supply a remainder below one third of
that margin at any containing section with $n\le800$.
\end{proposition}

\section{Kill and pad audit}

\begin{center}
\small
\begin{tabular}{@{}p{0.31\linewidth}p{0.58\linewidth}@{}}
\toprule
Binding item & R496 disposition\\
\midrule
Galerkin as operator & blocked; section values are obstruction
diagnostics only\\
Unbooked $[0,x_0]$ & blocked; every row uses its actual $\kappa_w$\\
HS without full $3\times$ charge & blocked; $r+2r=3r$ is reported\\
Frequency re-entry & blocked after the five-function G1 seal\\
Prime support convention & r495-A2 zero-extension battery passes
$6/6$ exactly\\
Prime section treated as finite rank & rejected; nonvanishing
off-block is measured explicitly\\
\bottomrule
\end{tabular}
\end{center}

\section{What remains open}

The r471 grid certificate and r477 two-dimensional certificate are
finite-subspace positives, not lower bounds for $\lambda_*(0.8)$.
This round does not decide whether another representation can prove
full-class positivity at $L=0.8$.  Any successor would need a
noncompact-translation argument, rather than the compact
Hilbert--Schmidt completion that closed r494.  Per the r492 budget,
this lane stops here.

\end{document}
